\subsubsection{অংশের অর্থ বোঝানো ট্যাগ: আধুনিক পরিচয়}
একটি কলেজের পেজে অনেক \texttt{div} থাকলে কোনটি মেনু আর কোনটি মূল লেখা, শুধু ট্যাগ দেখে বোঝা যায় না। \textbf{অর্থবোধক এলিমেন্ট (Semantic Element)} অংশের ভূমিকা প্রকাশ করে। এতে কোডের গঠন বোঝা এবং স্ক্রিন রিডার দিয়ে অংশ খুঁজে পাওয়া সহজ হয়। এগুলো নিজেরাই পেজের রং বা কলামের বিন্যাস নির্ধারণ করে না।

\begin{tabular}{|p{.22\linewidth}|p{.68\linewidth}|}
\hline
\textbf{এলিমেন্ট} & \textbf{অর্থ ও কলেজের পেজে প্রয়োগ} \\
\hline
\texttt{header} & পেজ বা বিভাগের পরিচিতি; কলেজের নাম ও লোগো \\
\hline
\texttt{nav} & গুরুত্বপূর্ণ নেভিগেশন লিংকের অংশ; প্রধান মেনু \\
\hline
\texttt{main} & পেজের প্রধান স্বতন্ত্র বিষয়বস্তু; নোটিশ বা পরিচিতি \\
\hline
\texttt{section} & বিষয়ভিত্তিক বিভাগ; ভর্তি-তথ্যের আলাদা অংশ, সাধারণত শিরোনামসহ \\
\hline
\texttt{article} & স্বয়ংসম্পূর্ণ লেখা; আলাদা করে পড়া যায় এমন নোটিশ \\
\hline
\texttt{aside} & মূল বিষয়ের সঙ্গে পরোক্ষভাবে সম্পর্কিত অংশ; সহায়ক তথ্য \\
\hline
\texttt{footer} & পেজ বা বিভাগের শেষের তথ্য; যোগাযোগ বা কপিরাইট \\
\hline
\end{tabular}

অংশের নির্দিষ্ট ভূমিকা থাকলে উপযুক্ত অর্থবোধক ট্যাগ ব্যবহার করা যায়। কেবল কয়েকটি এলিমেন্ট দলবদ্ধ করে CSS প্রয়োগের জন্য সাধারণ ধারক div উপযোগী।

